% Main-manuscript tables, in their original numerical order.
% Kept separate from SI tables and from the journal layout.

\begin{table}[H]
\centering
\caption{\textbf{Hidden Markov model with jumps (HMM-WJ) and
    single-index model (SIM) composition hyperparameters.} Values
    used for the main SPY analysis, explored ranges, and locations of
    sensitivity checks. For the primary six-method SIM comparison, we set
    $\epsilon = 0$ instead of the SPY value. The separate jump-enabled
    comparison transferred the listed SPY settings to enabled models. The cross-asset
    validation used per-ticker $(\epsilon^{*}, \lambda^{*})$ values
    for NVDA, JNJ, and JPM (Table~\ref{tab:cross_asset}).}
\label{tab:hyperparams}
\small
\begin{tabular}{@{}llll@{}}
\toprule
\textbf{Symbol} & \textbf{Description} & \textbf{SPY value} & \textbf{Explored range} \\ \midrule
\multicolumn{4}{l}{\textit{State partition and emission}} \\[2pt]
$N$            & Laplace-quantile states                      & $100$     & $\{30,\dots,350\}$, Online Appendix~\ref{sec:supp-sensitivity} \\
$\nu$          & Student-$t$ emission degrees of freedom      & $5$       & $\{3,\dots,30,\infty\}$ \\[4pt]
\multicolumn{4}{l}{\textit{Jump-duration override}} \\[2pt]
$\epsilon$     & jump probability per step                    & $10^{-4}$ & $[10^{-4}, 2.5\times10^{-2}]$, Eq.~\eqref{eq:grid_search} \\
$\lambda$      & mean jump-episode duration (steps)           & $90$      & $[10, 160]$, Eq.~\eqref{eq:grid_search} \\
$p_{\rm neg}$  & probability a jump targets $\mathcal{S}_{-}$ & $0.52$    & fixed \\
$N_{\rm tail}$ & quantile states per tail set                 & $5$       & fixed (at $N = 100$) \\
$w_\kappa$     & kurtosis-penalty weight, Eq.~\eqref{eq:grid_search} & $0.20$ & fixed \\
$L$            & autocorrelation lag window, Eq.~\eqref{eq:grid_search} & $25$ days & fixed \\[4pt]
\multicolumn{4}{l}{\textit{SIM composition}} \\[2pt]
$f$                        & idiosyncratic variance floor & $0.10$ & $\{0.05,\dots,0.30\}$ \\
$R^2_{\mathrm{preserve}}$  & branch-selection threshold   & $0.80$ & $\{0.70,\dots,0.90\}$ \\
\bottomrule
\end{tabular}
\end{table}

\begin{table}[tbp]
\centering
\caption{\textbf{In-sample and out-of-sample model comparison for SPY.}
Results summarize $1{,}000$ simulated paths. Values in parentheses are
standard errors, and bold entries are best among the fitted generative
models within each panel; the empirical bootstrap is excluded from
this highlighting. Metric
definitions and standard-error calculations are given in Online
Appendix~\ref{sec:supp-validation}.}
\label{tab:model_comparison}
\footnotesize
\setlength{\tabcolsep}{1.8pt}
\renewcommand{\arraystretch}{1.08}
\resizebox{\textwidth}{!}{%
\begin{tabular}{@{}lccccccc@{}}
\toprule
\textbf{Model} & \textbf{KS (\%)} & \textbf{AD (\%)} & $\boldsymbol{\kappa}$ &
\textbf{ACF-MAE} & \textbf{Cov. (\%)} & $\boldsymbol{W_1}$ & \textbf{Hellinger} \\
\midrule
\multicolumn{8}{@{}l}{\textit{In sample: $2014$--$2024$ ($2{,}766$ days; observed $\kappa=7.715$)}} \\
Bootstrap & 100.0 ($<$0.1) & 99.7 (0.2) & 7.6 (0.08) & 0.060 ($<$0.001) & 100.0 ($<$0.1) & 0.062 (0.001) & 0.042 ($<$0.001) \\
Gaussian  & 0.0 ($<$0.1) & 0.0 ($<$0.1) & $-$0.0 ($<$0.01) & 0.060 ($<$0.001) & 13.1 (0.1) & 0.399 (0.001) & 0.148 ($<$0.001) \\
Laplace   & 44.0 (1.6) & 43.3 (1.6) & 3.0 (0.02) & 0.060 ($<$0.001) & 37.4 (1.5) & 0.138 (0.001) & \textbf{0.072} ($<$0.001) \\
GARCH(1,1) & 5.5 (0.7) & 1.9 (0.4) & 8.2 (0.37) & \textbf{0.031} (0.001) & 29.3 (1.4) & 0.304 (0.006) & 0.100 (0.001) \\
GRU       & 0.6 (0.2) & 0.2 (0.1) & 5.7 (0.16) & 0.036 ($<$0.001) & 17.2 (0.5) & 0.421 (0.003) & 0.134 ($<$0.001) \\
HSMM      & 82.0 (1.2) & 42.5 (1.6) & 4.8 (0.08) & 0.059 ($<$0.001) & 68.7 (1.1) & 0.176 (0.001) & 0.113 ($<$0.001) \\
HMM-NJ    & \textbf{99.3} (0.3) & \textbf{98.2} (0.4) & 8.1 (0.14) & 0.059 ($<$0.001) & \textbf{100.0} ($<$0.1) & \textbf{0.081} (0.001) & \textbf{0.072} ($<$0.001) \\
HMM-WJ    & 98.3 (0.4) & 94.8 (0.7) & \textbf{7.5} (0.09) & 0.053 ($<$0.001) & \textbf{100.0} ($<$0.1) & 0.097 (0.001) & 0.074 ($<$0.001) \\[3pt]
\midrule
\multicolumn{8}{@{}l}{\textit{Out of sample: $2025$ ($249$ days; observed $\kappa=6.867$)}} \\
Bootstrap & 97.4 (0.5) & 98.5 (0.4) & 6.0 (0.18) & 0.043 ($<$0.001) & \textbf{100.0} ($<$0.1) & 0.232 (0.002) & 0.205 (0.001) \\
Gaussian  & 62.2 (1.5) & 46.3 (1.6) & $-$0.0 ($<$0.01) & 0.043 ($<$0.001) & 36.4 (1.9) & 0.452 (0.002) & 0.235 (0.001) \\
Laplace   & 88.0 (1.0) & 92.9 (0.8) & 2.7 (0.05) & 0.043 ($<$0.001) & 74.7 (1.0) & 0.263 (0.002) & 0.211 (0.001) \\
GARCH(1,1) & 80.3 (1.3) & 72.0 (1.4) & 1.6 (0.07) & \textbf{0.026} ($<$0.001) & 96.0 (1.3) & 0.507 (0.018) & 0.232 (0.001) \\
GRU       & 71.8 (1.4) & 44.8 (1.6) & 2.9 (0.10) & 0.033 ($<$0.001) & 77.8 (2.9) & 0.488 (0.005) & 0.240 (0.001) \\
HSMM      & 96.2 (0.6) & \textbf{96.7} (0.6) & 4.1 (0.13) & 0.042 ($<$0.001) & \textbf{100.0} (0.2) & 0.287 (0.002) & 0.239 (0.001) \\
HMM-NJ    & \textbf{96.6} (0.6) & 96.1 (0.6) & \textbf{6.2} (0.15) & 0.041 ($<$0.001) & \textbf{100.0} ($<$0.1) & \textbf{0.259} (0.003) & \textbf{0.207} (0.001) \\
HMM-WJ    & 95.4 (0.7) & 95.3 (0.7) & \textbf{6.2} (0.15) & 0.040 ($<$0.001) & \textbf{100.0} ($<$0.1) & 0.275 (0.005) & 0.210 (0.001) \\
\bottomrule
\end{tabular}}
\vspace{3pt}

\parbox{0.98\textwidth}{\scriptsize
KS and AD report pass rates at $\alpha=0.05$; coverage is the fraction
of empirical quantiles inside the synthetic $90\%$ envelope.
$\kappa$ is simulated excess kurtosis. ACF-MAE is computed for
$|G_t|$ through lag $252$. Larger values are preferred for KS, AD,
and coverage; smaller values are preferred for ACF-MAE, $W_1$, and
Hellinger distance; $\kappa$ is compared with the observed value.}
\end{table}


\begin{table}[H]
  \centering
  \caption{\textbf{Aggregate scorecard across $423$ non-market tickers and
    $100$ replications per ticker.} Errors, distances, and tail metrics
    are medians per method. Kolmogorov--Smirnov (KS) and
    Anderson--Darling (AD) pass fractions are pooled across replications.
    $|\hat\alpha - \alpha|$, $|\hat\beta - \beta|$, and
    $|\hat R^2 - R^2_{\real}|$: absolute deviation between ordinary
    least-squares (OLS) recovery
    and calibration (intercept error has units of $\mathrm{yr}^{-1}$).
    KS pass and AD pass: fraction of
    replications whose two-sample test against the real marginal exceeds
    $p = 0.05$. $W_1$: Wasserstein-1 distance to the real marginal.
    $\kappa$: sample excess kurtosis. Hill: tail-index estimate on the
    upper $5\%$ of $|g|$. SIM denotes the single-index model, and GARCH
    denotes generalized autoregressive conditional heteroskedasticity.}
  \label{tab:aggregate}
  \small
  \resizebox{\textwidth}{!}{\begin{tabular}{lrrrrrrrr}
\toprule
Composer & $|\hat\alpha-\alpha|$ & $|\hat\beta-\beta|$ & $|\hat R^2 - R^2_{\real}|$ & KS pass (\%) & AD pass (\%) & $W_1$ & $\kappa$ & Hill \\
\midrule
Naive & 0.002 & 0.022 & 0.087 & 7.1 & 3.5 & 0.700 & 7.767 & 0.369 \\
Gaussian SIM & 0.048 & 0.017 & 0.008 & 0.6 & 0.5 & 0.682 & 1.427 & 0.217 \\
Hybrid & 0.002 & 0.018 & 0.021 & 73.6 & 64.6 & 0.249 & 7.165 & 0.365 \\
JumpHMM-on-residuals & 0.049 & 0.018 & 0.015 & 75.8 & 70.7 & 0.232 & 7.262 & 0.365 \\
Block bootstrap & 0.042 & 0.017 & 0.020 & 73.3 & 66.8 & 0.232 & 6.675 & 0.330 \\
GARCH(1,1)-$t$ & 0.047 & 0.017 & 0.036 & 67.4 & 61.4 & 0.267 & 6.069 & 0.317 \\
\bottomrule
\end{tabular}
}
\end{table}

\begin{table}[H]
  \centering
  \caption{\textbf{Frozen-parameter multi-asset evaluation on the
    $2025$ holdout.} Results cover $416$ non-market tickers and $100$
    replications per ticker; generalized autoregressive conditional
    heteroskedasticity (GARCH) covers $386$ tickers. Pass rates
    are first averaged within ticker and then across tickers.
    The \textit{matched in-sample (IS)} columns compare each synthetic path with a randomly
    selected contiguous $249$-day block from the $2014$--$2024$
    training history, matching the holdout sample length. $W_1$ and
    the variance ratio are medians across ticker-level medians; the
    variance ratio is synthetic variance divided by observed $2025$
    variance. The hybrid correction used the training-period market
    variance and each simulated asset draw's variance.}
  \label{tab:oos_composition}
  \small
  \begin{tabular}{lrrrrrr}
\toprule
 & \multicolumn{2}{c}{KS pass (\%)} & \multicolumn{2}{c}{AD pass (\%)} & & \\
\cmidrule(lr){2-3} \cmidrule(lr){4-5}
Composer & matched IS & 2025 & matched IS & 2025 & $W_1$ (2025) & variance ratio \\
\midrule
Naive & 69.4 & 83.5 & 54.8 & 71.6 & 0.877 & 1.29 \\
Gaussian SIM & 57.4 & 73.9 & 44.9 & 64.6 & 0.866 & 0.97 \\
Hybrid & 78.8 & 87.0 & 66.2 & 79.7 & 0.716 & 0.97 \\
JumpHMM-on-residuals & 77.9 & 86.2 & 65.6 & 80.1 & 0.708 & 0.98 \\
Block bootstrap & 76.3 & 84.0 & 64.0 & 76.5 & 0.715 & 0.90 \\
GARCH(1,1)-$t$ & 74.5 & 82.6 & 62.6 & 75.0 & 0.733 & 0.91 \\
\bottomrule
\end{tabular}

\end{table}

\begin{table}[H]
  \centering
  \caption{\textbf{Out-of-sample per-ticker one-day left-tail Value-at-Risk
    coverage across the six composition methods.} VaR thresholds at
    coverage level $\alpha$ were estimated from frozen-parameter
    synthetic paths and exceedances were counted against the real
    $2025$ growth-rate history. The hybrid correction used the same
    training-period market variance convention as
    Table~\ref{tab:oos_composition}. \textit{rate}: mean of the per-ticker
    mean exceedance rates (nominal: $1-\alpha$). Cross-ticker standard
    deviation (SD) is reported in percentage points.
    \textit{Kupiec pass}: mean within-ticker fraction of replications
    whose unconditional-coverage test $p$-value exceeds $0.05$.}
  \label{tab:var}
  \small
  \setlength{\tabcolsep}{5pt}
  \begin{tabular}{lrrrrrr}
\toprule
 & \multicolumn{3}{c}{$\alpha = 0.95$} & \multicolumn{3}{c}{$\alpha = 0.99$} \\
\cmidrule(lr){2-4} \cmidrule(lr){5-7}
Composer & rate (\%) & SD (pp) & Kupiec pass (\%) & rate (\%) & SD (pp) & Kupiec pass (\%) \\
\midrule
Naive & 4.37 & 2.15 & 65.8 & 1.21 & 0.76 & 78.9 \\
Gaussian SIM & 4.95 & 2.26 & 70.0 & 2.03 & 1.17 & 72.1 \\
Hybrid & 5.97 & 2.43 & 66.3 & 1.67 & 0.93 & 73.4 \\
JumpHMM-on-residuals & 6.07 & 2.59 & 67.0 & 1.67 & 0.95 & 74.5 \\
Block bootstrap & 6.13 & 2.56 & 63.7 & 1.84 & 1.01 & 73.4 \\
GARCH(1,1)-$t$ & 5.79 & 2.47 & 62.4 & 1.80 & 1.02 & 72.0 \\
\bottomrule
\end{tabular}

\end{table}

\begin{table}[H]
  \centering
  \caption{\textbf{Jump settings under naive and corrected multi-asset
    composition.} Each setting used $1{,}000$ replications per asset
    across four seeds, with identical asset and market draws for the
    two composition methods. SPY was simulated in both windows,
    and the correction used each simulated market path's variance
    and each asset draw's variance. Entries are equal-weight
    cross-ticker means of per-path metrics.
    Kolmogorov--Smirnov (KS) and Anderson--Darling (AD) report
    non-rejection rates at $\alpha=0.05$. ACF$_{25}$ and ACF$_{60}$
    are mean absolute errors in the absolute-return autocorrelation
    function through lags $25$ and $60$; lower values indicate better
    temporal fit. $V_{\gen}=\Var(g_i)/\sigma_{\gen,i}^2$ compares
    composed variance with the active generator's variance, rather
    than observed variance.}
  \label{tab:jump_ablation}
  \footnotesize
  \setlength{\tabcolsep}{4pt}
  % Generated by 11c-Jump-Ablation-Tables.py; do not edit numbers manually.
\begin{tabular}{@{}llrrrrr@{}}
\toprule
Jump setting & Method & KS (\%) & AD (\%) & ACF$_{25}$ & ACF$_{60}$ & $V_{\gen}$ \\
\midrule
\multicolumn{7}{@{}l}{\textit{Training: $2014$--$2024$, $423$ assets, $2{,}766$ days}} \\[2pt]
Off & Naive & 6.93 & 3.38 & 0.13585 & 0.09492 & 1.33273 \\
Off & Corrected & 73.94 & 68.54 & 0.13569 & 0.09485 & 1.00011 \\
Market only & Naive & 6.28 & 3.18 & 0.13197 & 0.09208 & 1.34938 \\
Market only & Corrected & 73.64 & 67.81 & 0.12844 & 0.08969 & 1.00032 \\
Market + assets & Naive & 5.25 & 2.56 & 0.11727 & 0.08264 & 1.33445 \\
Market + assets & Corrected & 61.68 & 53.33 & 0.11581 & 0.08154 & 1.00018 \\
\midrule
\multicolumn{7}{@{}l}{\textit{Holdout: $2025$, $416$ assets, $249$ days}} \\[2pt]
Off & Naive & 82.38 & 70.14 & 0.07322 & 0.06683 & 1.36081 \\
Off & Corrected & 85.92 & 78.76 & 0.07330 & 0.06687 & 1.00041 \\
Market only & Naive & 81.76 & 69.60 & 0.07347 & 0.06698 & 1.36638 \\
Market only & Corrected & 85.88 & 78.72 & 0.07525 & 0.06796 & 1.00045 \\
Market + assets & Naive & 80.27 & 68.19 & 0.07672 & 0.06885 & 1.36032 \\
Market + assets & Corrected & 84.43 & 77.21 & 0.07832 & 0.06973 & 1.00039 \\
\bottomrule
\end{tabular}

\end{table}
